%=============================================================
%  TIERED BIBLIOGRAPHY
%  Quantum Mechanics & Quantum Chemistry Programme
%  One section per learning tier (HS → PhD) + Full Catalogue
%=============================================================
\documentclass[11pt,a4paper]{article}

\usepackage[T1]{fontenc}
\usepackage[utf8]{inputenc}
\usepackage{fix-cm}
\usepackage[margin=2.5cm]{geometry}
\usepackage{amsmath,amssymb}
\usepackage{booktabs}
\usepackage{array}
\usepackage{tabularx}
\usepackage{enumitem}
\usepackage{sectsty}
\usepackage{fancyhdr}
\usepackage[dvipsnames,svgnames,table]{xcolor}
\usepackage{hyperref}
\usepackage{parskip}
\usepackage{tcolorbox}
\tcbuselibrary{skins,breakable}
\usepackage{mdframed}

%--- Colours -------------------------------------------------
\definecolor{darkblue}{RGB}{10,40,90}
\definecolor{midblue}{RGB}{30,80,160}
\definecolor{lightblue}{RGB}{220,233,255}
\definecolor{accentgold}{RGB}{180,140,20}
\definecolor{lightgold}{RGB}{255,248,220}
\definecolor{darkgray}{RGB}{60,60,60}
% Tier colours
\definecolor{t1green}{RGB}{30,120,50}
\definecolor{t1light}{RGB}{212,242,220}
\definecolor{t2blue}{RGB}{25,90,170}
\definecolor{t2light}{RGB}{215,232,255}
\definecolor{t3purple}{RGB}{110,35,160}
\definecolor{t3light}{RGB}{238,218,255}
\definecolor{t4orange}{RGB}{165,70,10}
\definecolor{t4light}{RGB}{255,232,208}
\definecolor{t5red}{RGB}{140,15,15}
\definecolor{t5light}{RGB}{255,215,215}
\definecolor{catgray}{RGB}{80,80,80}
\definecolor{catlight}{RGB}{240,240,240}

%--- Section fonts -------------------------------------------
\sectionfont{\Large\bfseries\color{darkblue}\sectionrule{0ex}{0pt}{-1ex}{1.5pt}}
\subsectionfont{\large\bfseries\color{midblue}}
\subsubsectionfont{\normalsize\bfseries\color{darkgray}}

%--- Hyperref ------------------------------------------------
\hypersetup{
  colorlinks=true, linkcolor=midblue, urlcolor=midblue,
  pdftitle={QM Programme --- Tiered Bibliography},
}

%--- Header / Footer -----------------------------------------
\pagestyle{fancy}
\fancyhf{}
\fancyhead[L]{\small\color{darkgray}\textit{QM \& QC Programme --- Tiered Bibliography}}
\fancyhead[R]{\small\color{darkgray}\thepage}
\fancyfoot[C]{\small\color{darkgray}\textit{Companion to the Pedagogical Support Material}}
\renewcommand{\headrulewidth}{0.4pt}

%--- tcolorbox styles ----------------------------------------
\tcbset{
  tierheader/.style={
    enhanced, colback=#1, colframe=#1!80!black,
    boxrule=0pt, arc=0pt, sharp corners,
    left=8pt, right=8pt, top=6pt, bottom=6pt,
    fontupper=\bfseries\large\color{white},
  },
  entrybox/.style={
    enhanced, breakable,
    colback=#1, colframe=#1!60!black,
    boxrule=0.5pt, sharp corners,
    left=8pt, right=8pt, top=5pt, bottom=5pt,
  },
  annotbox/.style={
    enhanced, breakable,
    colback=lightgold, colframe=accentgold,
    boxrule=0.6pt, sharp corners,
    left=5pt, right=5pt, top=3pt, bottom=3pt,
    fontupper=\small\itshape\color{darkgray},
  },
  keybox/.style={
    enhanced,
    colback=lightblue, colframe=midblue,
    boxrule=0.5pt, rounded corners,
    left=5pt, right=5pt, top=4pt, bottom=4pt,
    fontupper=\small,
  },
  catbox/.style={
    enhanced, breakable,
    colback=catlight, colframe=catgray,
    boxrule=0.6pt, sharp corners,
    left=6pt, right=6pt, top=5pt, bottom=5pt,
  },
}

% Tier badge command
\newcommand{\tierbadge}[2]{%
  \colorbox{#1}{\small\bfseries\color{white}\strut\ #2\ }%
}

% Star rating
\newcommand{\essential}{$\bigstar\bigstar\bigstar$}
\newcommand{\recommended}{$\bigstar\bigstar$}
\newcommand{\useful}{$\bigstar$}

% Entry formatting
\newenvironment{bibentry}[1]{%
  \begin{tcolorbox}[entrybox=#1, top=4pt, bottom=4pt]%
}{%
  \end{tcolorbox}\smallskip%
}

%=============================================================
\begin{document}
%=============================================================

%--- Title Page ----------------------------------------------
\begin{titlepage}
\pagecolor{darkblue}
\color{white}
\vspace*{3cm}
\begin{center}
  {\fontsize{28}{34}\selectfont\bfseries Tiered Bibliography}\\[0.5cm]
  {\fontsize{17}{22}\selectfont Quantum Mechanics \& Quantum Chemistry Programme}\\[0.6cm]
  {\color{accentgold}\rule{10cm}{2pt}}\\[0.6cm]
  {\large\itshape Reading Lists Stratified by Learning Tier}\\[2.5cm]
  \begin{tabular}{@{}cl@{}}
    \colorbox{t1green}{\color{white}\bfseries Tier 1} & High School --- Popular Science \& First Introductions\\[4pt]
    \colorbox{t2blue}{\color{white}\bfseries Tier 2}  & Beginning Undergraduate --- Core Textbooks\\[4pt]
    \colorbox{t3purple}{\color{white}\bfseries Tier 3} & Advanced Undergraduate --- Rigorous \& Bridge Texts\\[4pt]
    \colorbox{t4orange}{\color{white}\bfseries Tier 4} & Master's --- Graduate Textbooks \& Key Papers\\[4pt]
    \colorbox{t5red}{\color{white}\bfseries Tier 5}   & PhD --- Monographs, Reviews \& Primary Literature\\[8pt]
    \colorbox{catgray}{\color{white}\bfseries Cat.~A--F} & Full Alphabetical Catalogue (all tiers)\\
  \end{tabular}\\[2.5cm]
  {\color{accentgold}\rule{10cm}{1pt}}\\[0.5cm]
  {\small\itshape Companion document to the Pedagogical Support Material.\\
  \textbf{Annotation key:}
  \essential\ Essential \quad
  \recommended\ Strongly recommended \quad
  \useful\ Useful supplement\\[4pt]
  \textbf{[P]} Primary paper \quad
  \textbf{[R]} Review \quad
  \textbf{[M]} Monograph \quad
  \textbf{[N]} Popular/narrative}
\end{center}
\vfill
\end{titlepage}
\nopagecolor

\tableofcontents
\newpage

%=============================================================
\section*{How to Use This Bibliography}
\addcontentsline{toc}{section}{How to Use This Bibliography}
%=============================================================

\begin{tcolorbox}[keybox]
This bibliography is a companion to the \textbf{Pedagogical Support Material} (PSM),
which describes the five-tier learning system and the 14-type assessment artifact set.
Each tier corresponds to a defined competency level; the reading lists here feed
directly into the PSM's learning objectives, concept questions, and research tracks.
\end{tcolorbox}

\medskip
\begin{center}
\begin{tabular}{clp{3.5cm}p{5.5cm}}
\toprule
\textbf{Tier} & \textbf{Label} & \textbf{PSM track} & \textbf{Reading strategy}\\
\midrule
1 & HS       & Undergrad & Popular books cover-to-cover first\\
2 & BegUG    & Undergrad & One textbook chapter per lecture\\
3 & AdvUG    & Undergrad & Textbook + primary paper abstracts\\
4 & MSc      & Graduate  & Primary papers as main reading\\
5 & PhD      & Graduate + Research & Monographs + arXiv preprints\\
\bottomrule
\end{tabular}
\end{center}

\medskip
\noindent The tiers are \textbf{cumulative}: all Tier~$k$ reading is recommended as
background for Tier~$k{+}1$. Each entry carries:
\begin{itemize}[leftmargin=1.8em,itemsep=1pt]
  \item A \textbf{star rating} (\essential\ / \recommended\ / \useful)
  \item A \textbf{type tag} ([P] [R] [M] [N])
  \item A \textbf{short annotation} linking the work to specific programme modules
\end{itemize}

\medskip
\noindent\textbf{Quantum Mechanics} entries are always primary; cross-references to
\textbf{Nonlinear Dynamics \& Chaos} (the logistic map course) are included at each
tier to illustrate how the same tier architecture transfers across disciplines.

\newpage

%#############################################################
\section{Tier 1 (HS) --- Popular Science \& First Introductions}
\label{sec:t1}
%#############################################################

\begin{tcolorbox}[tierheader=t1green]
Tier 1 --- High School (HS)
\end{tcolorbox}

\medskip
\begin{tcolorbox}[keybox]
\textbf{Goal:} Build vivid physical intuition and historical awareness before
encountering any formalism. All books in this tier require only high-school mathematics.
Read these \emph{before} any formal course. They provide the mental pictures that make
the formalism of Tiers~2--5 meaningful.
\end{tcolorbox}

%------------------------------------------------------------
\subsection{Quantum Mechanics}
%------------------------------------------------------------

\begin{bibentry}{t1light}
\textbf{Feynman, R.~P.}\\
\textit{QED: The Strange Theory of Light and Matter}\\
Princeton University Press, 1985.\\[4pt]
\essential\ \ [\textbf{N}]\\[4pt]
\begin{tcolorbox}[annotbox]
The finest popular account of quantum amplitudes and interference. Feynman derives
the Born rule from ``arrows'' that rotate and add --- precisely the Tier~1 picture
of the bra-ket formalism. The ``photon bouncing off glass'' chapters are a perfect
pre-reading for Module~I.2, Lectures~L01--L02. Read before any formal course.
\end{tcolorbox}
\end{bibentry}

\begin{bibentry}{t1light}
\textbf{Feynman, R.~P., Leighton, R.~B.\ \& Sands, M.}\\
\textit{The Feynman Lectures on Physics, Vol.~III: Quantum Mechanics}\\
Addison-Wesley, 1965. Freely available at \texttt{feynmanlectures.caltech.edu}\\[4pt]
\essential\ \ [\textbf{N/pedagogical}]\\[4pt]
\begin{tcolorbox}[annotbox]
Volume~III opens with two-slit interference and amplitudes; the first four chapters
are fully accessible at Tier~1 and map directly onto the Tier~1 engagement activities
for Module~I.2 Lectures L01 and L02. An essential free resource.
\end{tcolorbox}
\end{bibentry}

\begin{bibentry}{t1light}
\textbf{Cox, B.\ \& Forshaw, J.}\\
\textit{The Quantum Universe: Everything That Can Happen Does Happen}\\
Da Capo Press, 2011.\\[4pt]
\recommended\ \ [\textbf{N}]\\[4pt]
\begin{tcolorbox}[annotbox]
Readable narrative covering superposition, entanglement, and the Pauli exclusion
principle with minimal formalism. Good companion for Module~I.2 Lectures L09--L10
(tensor products, spin). Can be read in a weekend.
\end{tcolorbox}
\end{bibentry}

\begin{bibentry}{t1light}
\textbf{Kumar, M.}\\
\textit{Quantum: Einstein, Bohr, and the Great Debate about the Nature of Reality}\\
Norton, 2008.\\[4pt]
\recommended\ \ [\textbf{N}]\\[4pt]
\begin{tcolorbox}[annotbox]
Historical narrative of the foundational debates about measurement and probability.
Essential context for understanding why Born's rule (1926) was controversial. Prepares
the historical framing for Module~I.2 Lectures L04 and L07.
\end{tcolorbox}
\end{bibentry}

\begin{bibentry}{t1light}
\textbf{Penrose, R.}\\
\textit{The Emperor's New Mind}, Chapters~5--6\\
Oxford University Press, 1989.\\[4pt]
\useful\ \ [\textbf{N}]\\[4pt]
\begin{tcolorbox}[annotbox]
Philosophical context for quantum mechanics. Chapters~5--6 discuss wave functions and
measurement in an accessible way. Useful philosophical background before Module~I.3
Lectures L05 and L07.
\end{tcolorbox}
\end{bibentry}

%------------------------------------------------------------
\subsection{Nonlinear Dynamics \& Chaos (Cross-reference)}
%------------------------------------------------------------

\begin{bibentry}{t1light}
\textbf{Gleick, J.}\\
\textit{Chaos: Making a New Science}\\
Penguin Books, 1987.\\[4pt]
\essential\ \ [\textbf{N}]\\[4pt]
\begin{tcolorbox}[annotbox]
Narrative history of May, Feigenbaum, and the logistic map's role in the chaos
revolution. No mathematics required. The exact Tier~1 equivalent of Feynman's QED
for nonlinear dynamics. Standard first reading for any dynamical systems course.
\end{tcolorbox}
\end{bibentry}

\begin{bibentry}{t1light}
\textbf{Stewart, I.}\\
\textit{Does God Play Dice? The New Mathematics of Chaos}\\
Blackwell, 1989.\\[4pt]
\recommended\ \ [\textbf{N}]\\[4pt]
\begin{tcolorbox}[annotbox]
Gentle narrative; introduces bifurcation and sensitive dependence at high-school level.
\end{tcolorbox}
\end{bibentry}

\newpage
%#############################################################
\section{Tier 2 (BegUG) --- Core Undergraduate Textbooks}
\label{sec:t2}
%#############################################################

\begin{tcolorbox}[tierheader=t2blue]
Tier 2 --- Beginning Undergraduate (BegUG)
\end{tcolorbox}

\medskip
\begin{tcolorbox}[keybox]
\textbf{Goal:} Develop working fluency in the notation and methods of quantum mechanics.
All books assume first-year linear algebra and calculus. Work through one chapter per
lecture, solving end-of-chapter exercises alongside SPSU problem sets from the PSM.
\end{tcolorbox}

%------------------------------------------------------------
\subsection{Quantum Mechanics}
%------------------------------------------------------------

\begin{bibentry}{t2light}
\textbf{Griffiths, D.~J.\ \& Schroeter, D.~F.}\\
\textit{Introduction to Quantum Mechanics} (3rd ed.)\\
Cambridge University Press, 2018.\\[4pt]
\essential\\[4pt]
\begin{tcolorbox}[annotbox]
The standard first course. Develops wave mechanics before bra-ket notation.
Chapters~1--3 (wavefunction, Schr\"odinger equation, infinite square well) parallel
Module~I.3 Lectures L02--L05; Chapter~3 introduces Hilbert space and Dirac notation,
paralleling Module~I.2 Lectures L01--L05. Excellent problem sets map to SPSU format.
\end{tcolorbox}
\end{bibentry}

\begin{bibentry}{t2light}
\textbf{Shankar, R.}\\
\textit{Principles of Quantum Mechanics} (2nd ed.)\\
Springer, 1994.\\[4pt]
\essential\\[4pt]
\begin{tcolorbox}[annotbox]
Begins with a self-contained linear algebra review in Chapter~1 --- ideal for
students bridging Tier~1 to Tier~2. More algebraic than Griffiths; a natural
bridge to Sakurai. Chapter~1 pairs with Module~I.1; Chapters~4--9 support
Module~I.2 Lectures L03--L08.
\end{tcolorbox}
\end{bibentry}

\begin{bibentry}{t2light}
\textbf{Townsend, J.~S.}\\
\textit{A Modern Approach to Quantum Mechanics} (2nd ed.)\\
University Science Books, 2012.\\[4pt]
\recommended\\[4pt]
\begin{tcolorbox}[annotbox]
Starts from spin-$\tfrac{1}{2}$ and Dirac notation from the very first chapter
(Sakurai-style). Excellent for building bra-ket fluency early. Module~I.2
can use this as the parallel text for Lectures L01--L04 and L10.
\end{tcolorbox}
\end{bibentry}

\begin{bibentry}{t2light}
\textbf{Binney, J.\ \& Skinner, D.}\\
\textit{The Physics of Quantum Mechanics}\\
Oxford University Press, 2014. Freely available at \texttt{www-thphys.physics.ox.ac.uk}\\[4pt]
\recommended\\[4pt]
\begin{tcolorbox}[annotbox]
British undergraduate style; uses bra-ket from Chapter~1. Excellent treatment of
measurement and operators. Good companion for Module~I.2 Lectures L03--L07.
Free access makes it especially attractive.
\end{tcolorbox}
\end{bibentry}

%------------------------------------------------------------
\subsection{Mathematics Support}
%------------------------------------------------------------

\begin{bibentry}{t2light}
\textbf{Strang, G.}\\
\textit{Introduction to Linear Algebra} (5th ed.)\\
Wellesley--Cambridge Press, 2016.\\[4pt]
\essential\\[4pt]
\begin{tcolorbox}[annotbox]
The prerequisite linear algebra text. Chapters~5--6 (eigenvalues, orthogonality) are
directly needed for quantum mechanics. Work through before or alongside Module~I.1.
\end{tcolorbox}
\end{bibentry}

\begin{bibentry}{t2light}
\textbf{Byron, F.~W.\ \& Fuller, R.~W.}\\
\textit{Mathematics of Classical and Quantum Physics}\\
Dover, 1992.\\[4pt]
\recommended\\[4pt]
\begin{tcolorbox}[annotbox]
Covers Hilbert spaces, operators, and distributions at a level accessible to advanced
Tier~2 students. A useful mathematical supplement for Module~I.2 Lectures L06 and L12.
\end{tcolorbox}
\end{bibentry}

%------------------------------------------------------------
\subsection{Quantum Chemistry Introduction}
%------------------------------------------------------------

\begin{bibentry}{t2light}
\textbf{Jensen, F.}\\
\textit{Introduction to Computational Chemistry} (3rd ed.)\\
Wiley, 2017.\\[4pt]
\recommended\\[4pt]
\begin{tcolorbox}[annotbox]
Accessible entry to computational quantum chemistry. Chapters~1--3 provide the Tier~2
background for Series~III. Students should read alongside Module~III.1.
\end{tcolorbox}
\end{bibentry}

%------------------------------------------------------------
\subsection{Nonlinear Dynamics Cross-reference}
%------------------------------------------------------------

\begin{bibentry}{t2light}
\textbf{Strogatz, S.~H.}\\
\textit{Nonlinear Dynamics and Chaos} (2nd ed.)\\
Westview Press, 2015.\\[4pt]
\essential\\[4pt]
\begin{tcolorbox}[annotbox]
The standard first course in nonlinear dynamics. Chapters~10--12 cover 1D maps,
period-doubling, and chaos with clear physical motivation. Exact Tier~2 analogue
of Griffiths.
\end{tcolorbox}
\end{bibentry}

\begin{bibentry}{t2light}
\textbf{Alligood, K.~T., Sauer, T.~D.\ \& Yorke, J.~A.}\\
\textit{Chaos: An Introduction to Dynamical Systems}\\
Springer, 1996.\\[4pt]
\recommended\\[4pt]
\begin{tcolorbox}[annotbox]
Chapter~1 develops the logistic map from scratch with full proofs accessible to
beginning undergraduates.
\end{tcolorbox}
\end{bibentry}

\newpage
%#############################################################
\section{Tier 3 (AdvUG) --- Rigorous Undergraduate \& Bridge Texts}
\label{sec:t3}
%#############################################################

\begin{tcolorbox}[tierheader=t3purple]
Tier 3 --- Advanced Undergraduate (AdvUG)
\end{tcolorbox}

\medskip
\begin{tcolorbox}[keybox]
\textbf{Goal:} Understand why the formalism is internally consistent and geometric.
Assumes confidence with analysis and linear algebra. Start reading primary papers
at the abstract and introduction level. Use alongside CPSU problem sets from the PSM.
\end{tcolorbox}

%------------------------------------------------------------
\subsection{Quantum Mechanics --- Core}
%------------------------------------------------------------

\begin{bibentry}{t3light}
\textbf{Sakurai, J.~J.\ \& Napolitano, J.}\\
\textit{Modern Quantum Mechanics} (3rd ed.)\\
Cambridge University Press, 2020.\\[4pt]
\essential\\[4pt]
\begin{tcolorbox}[annotbox]
The primary text for Module~I.2. Begins immediately with bra-ket notation and the
Stern--Gerlach experiment. Chapters~1--3 (state vectors, dynamics, angular momentum)
are the core of Series~I. The Sakurai ``operational'' approach is one of the two tracks
in the Module~I.2 dual pedagogy (Sakurai + Dirac). Essential reading at Tier~3.
\end{tcolorbox}
\end{bibentry}

\begin{bibentry}{t3light}
\textbf{Dirac, P.~A.~M.}\\
\textit{The Principles of Quantum Mechanics} (4th ed.)\\
Oxford University Press, 1958 (reprinted 1981).\\[4pt]
\essential\ \ [\textbf{P/pedagogical}]\\[4pt]
\begin{tcolorbox}[annotbox]
The original source for bra-ket notation. Not a modern textbook but an indispensable
primary reference. Chapter~I (``The Principle of Superposition'') and Chapter~II
(``Dynamical Variables and Observables'') should be read in full alongside
Module~I.2. This is the ``Dirac track'' of the dual pedagogy.
\end{tcolorbox}
\end{bibentry}

\begin{bibentry}{t3light}
\textbf{Isham, C.~J.}\\
\textit{Lectures on Quantum Theory: Mathematical and Structural Foundations}\\
Imperial College Press, 1995.\\[4pt]
\recommended\\[4pt]
\begin{tcolorbox}[annotbox]
Bridges undergraduate and graduate levels. Rigorous treatment of Hilbert spaces,
operators, and measurement without full functional analysis. Ideal for Tier~3
students aiming at Tier~4; directly supports Module~I.2 Lectures L11--L12.
\end{tcolorbox}
\end{bibentry}

\begin{bibentry}{t3light}
\textbf{Hannabuss, K.}\\
\textit{An Introduction to Quantum Theory}\\
Oxford University Press, 1997.\\[4pt]
\recommended\\[4pt]
\begin{tcolorbox}[annotbox]
Mathematically precise undergraduate text. Good treatment of spectral theory,
uncertainty, and density operators without measure theory. Supports Module~I.2
Lectures L04, L07, and L11.
\end{tcolorbox}
\end{bibentry}

%------------------------------------------------------------
\subsection{Primary Papers --- First Encounters}
%------------------------------------------------------------

\begin{bibentry}{t3light}
\textbf{Born, M.}\\
\textit{Zur Quantenmechanik der Sto\ss vorg\"ange}\\
Z.~Phys.~\textbf{37}, 863--867 (1926). English translation in Wheeler \& Zurek,
\textit{Quantum Theory and Measurement} (Princeton, 1983).\\[4pt]
\recommended\ \ [\textbf{P}]\\[4pt]
\begin{tcolorbox}[annotbox]
The Born rule. Two pages; the probability interpretation of $|\psi|^2$. Should be
read (in translation) by all Tier~3 students. Directly relevant to Module~I.2
Lecture L01 (``why probabilities are squared amplitudes'') and the WRQ prompt
``How did Born's interpretation change the meaning of the wavefunction?''
\end{tcolorbox}
\end{bibentry}

\begin{bibentry}{t3light}
\textbf{von Neumann, J.}\\
\textit{Mathematische Grundlagen der Quantenmechanik}\\
Springer, 1932. English: \textit{Mathematical Foundations of Quantum Mechanics},
Princeton University Press, 1955.\\[4pt]
\recommended\ \ [\textbf{M}]\\[4pt]
\begin{tcolorbox}[annotbox]
The mathematical foundation of quantum mechanics. Chapters~I--II (abstract Hilbert
space, operators) are accessible at Tier~3; Chapters~III--VI require Tier~5
preparation. Supports Module~I.2 Lecture L12 and Module~I.3 Lecture L07.
\end{tcolorbox}
\end{bibentry}

\begin{bibentry}{t3light}
\textbf{Wigner, E.~P.}\\
\textit{On unitary representations of the inhomogeneous Lorentz group}\\
Ann.~Math.~\textbf{40}, 149--204 (1939).\\[4pt]
\useful\ \ [\textbf{P}]\\[4pt]
\begin{tcolorbox}[annotbox]
Original source for Wigner's theorem and unitary/antiunitary symmetries. Abstract is
accessible at Tier~3; full proof requires Tier~5 preparation. Relevant to Module~I.2
Lectures L01 and L03 (``symmetries of physical states'').
\end{tcolorbox}
\end{bibentry}

%------------------------------------------------------------
\subsection{Nonlinear Dynamics Cross-reference}
%------------------------------------------------------------

\begin{bibentry}{t3light}
\textbf{Devaney, R.~L.}\\
\textit{An Introduction to Chaotic Dynamical Systems} (2nd ed.)\\
Addison-Wesley, 1989.\\[4pt]
\essential\\[4pt]
\begin{tcolorbox}[annotbox]
Mathematically precise. Introduces Devaney's three conditions for chaos. Bridge text
to graduate level; exact Tier~3 analogue of Sakurai for QM.
\end{tcolorbox}
\end{bibentry}

\begin{bibentry}{t3light}
\textbf{May, R.~M.}\\
\textit{Simple mathematical models with very complicated dynamics}\\
Nature \textbf{261}, 459--467 (1976).\\[4pt]
\essential\ \ [\textbf{P}]\\[4pt]
\begin{tcolorbox}[annotbox]
The landmark logistic map paper. Eight pages; entirely accessible at Tier~3.
The QM analogue for this tier is Born (1926). Should be read at Tier~3 in the
nonlinear dynamics course as the first primary paper encounter.
\end{tcolorbox}
\end{bibentry}

\newpage
%#############################################################
\section{Tier 4 (MSc) --- Graduate Textbooks \& Key Papers}
\label{sec:t4}
%#############################################################

\begin{tcolorbox}[tierheader=t4orange]
Tier 4 --- Master's (MSc)
\end{tcolorbox}

\medskip
\begin{tcolorbox}[keybox]
\textbf{Goal:} Command the full formalism; reproduce key derivations; situate results
in the broader literature. Primary papers become the main reading; textbooks serve
as reference. Functional analysis background assumed. Use alongside CPSG, SPJG,
and WRQ artifact types from the PSM.
\end{tcolorbox}

%------------------------------------------------------------
\subsection{Quantum Mechanics --- Graduate Texts}
%------------------------------------------------------------

\begin{bibentry}{t4light}
\textbf{Weinberg, S.}\\
\textit{Lectures on Quantum Mechanics} (2nd ed.)\\
Cambridge University Press, 2015.\\[4pt]
\essential\\[4pt]
\begin{tcolorbox}[annotbox]
Covers foundations (postulates, measurement, entanglement), symmetries, and
perturbation theory at graduate level. Chapter~3 (``Mathematical Tools'') is the
cleanest graduate treatment of the bra-ket formalism including rigged Hilbert spaces.
Supports Module~I.2 Lectures L11--L12 and Module~I.3 Lectures L01--L04.
\end{tcolorbox}
\end{bibentry}

\begin{bibentry}{t4light}
\textbf{Galindo, A.\ \& Pascual, P.}\\
\textit{Quantum Mechanics} (2 vols.)\\
Springer, 1990--1991.\\[4pt]
\essential\\[4pt]
\begin{tcolorbox}[annotbox]
The most mathematically complete standard graduate text. Volume~I develops Hilbert
space rigorously; Volume~II covers scattering, relativistic QM, and second
quantisation. Recommended for Tier~4 students aiming at mathematical physics.
Supports all of Series~I and II.
\end{tcolorbox}
\end{bibentry}

\begin{bibentry}{t4light}
\textbf{Ballentine, L.~E.}\\
\textit{Quantum Mechanics: A Modern Development} (2nd ed.)\\
World Scientific, 2014.\\[4pt]
\recommended\\[4pt]
\begin{tcolorbox}[annotbox]
Statistical interpretation; detailed treatment of measurement, density operators,
and the ensemble interpretation. Excellent for the foundational content of Module~I.2
Lectures L07 and L11, and Module~I.3 Lecture L10.
\end{tcolorbox}
\end{bibentry}

\begin{bibentry}{t4light}
\textbf{Cohen-Tannoudji, C., Diu, B.\ \& Lalo\"e, F.}\\
\textit{Quantum Mechanics} (2 vols.)\\
Wiley, 2019 (new edition).\\[4pt]
\recommended\\[4pt]
\begin{tcolorbox}[annotbox]
Encyclopaedic reference. The ``Complements'' to each chapter are especially valuable
for Tier~4 students who need worked examples alongside theory. Supports
Modules~II.3 (angular momentum) and~II.6 (spectroscopy).
\end{tcolorbox}
\end{bibentry}

%------------------------------------------------------------
\subsection{Many-Body Theory \& Advanced Topics}
%------------------------------------------------------------

\begin{bibentry}{t4light}
\textbf{Fetter, A.~L.\ \& Walecka, J.~D.}\\
\textit{Quantum Theory of Many-Particle Systems}\\
McGraw-Hill, 1971 (Dover reprint 2003).\\[4pt]
\essential\\[4pt]
\begin{tcolorbox}[annotbox]
The standard graduate text for Series~II Modules~II.4--II.5. Second quantisation,
Green's functions, and Feynman diagrams developed from scratch. Chapter~1 pairs
with Module~II.4 Lecture L01.
\end{tcolorbox}
\end{bibentry}

\begin{bibentry}{t4light}
\textbf{Bruus, H.\ \& Flensberg, K.}\\
\textit{Introduction to Many-Body Quantum Theory in Condensed Matter Physics}\\
Oxford University Press, 2004. Freely available at \texttt{community.dur.ac.uk}\\[4pt]
\essential\\[4pt]
\begin{tcolorbox}[annotbox]
More accessible than Fetter--Walecka. Excellent treatment of second quantisation,
diagrammatics, and Green's functions. Primary reading for Module~II.4.
Free access is a strong advantage.
\end{tcolorbox}
\end{bibentry}

\begin{bibentry}{t4light}
\textbf{Stefanucci, G.\ \& van Leeuwen, R.}\\
\textit{Nonequilibrium Many-Body Theory of Quantum Systems}\\
Cambridge University Press, 2013.\\[4pt]
\essential\\[4pt]
\begin{tcolorbox}[annotbox]
The definitive text for Module~II.5 (NEGF). Develops the Keldysh formalism
rigorously from first principles. Chapters~4--6 are the primary reading for
Module~II.5 Lectures L02--L06.
\end{tcolorbox}
\end{bibentry}

%------------------------------------------------------------
\subsection{Quantum Chemistry (Series III)}
%------------------------------------------------------------

\begin{bibentry}{t4light}
\textbf{Szabo, A.\ \& Ostlund, N.~S.}\\
\textit{Modern Quantum Chemistry: Introduction to Advanced Electronic Structure Theory}\\
McGraw-Hill, 1989 (Dover reprint 1996).\\[4pt]
\essential\\[4pt]
\begin{tcolorbox}[annotbox]
The standard entry-level graduate text for Series~III. Covers HF theory, CI, MP2,
and basic DFT concisely with full derivations. Chapters~1--3 support
Module~III.1; Chapters~4--6 support Module~III.2 Lectures L05--L11.
\end{tcolorbox}
\end{bibentry}

\begin{bibentry}{t4light}
\textbf{Helgaker, T., J\o rgensen, P.\ \& Olsen, J.}\\
\textit{Molecular Electronic-Structure Theory}\\
Wiley, 2000.\\[4pt]
\essential\ \ [\textbf{M}]\\[4pt]
\begin{tcolorbox}[annotbox]
The definitive graduate reference for Modules~III.2--III.3. 900 pages covering
basis sets, HF, CC, MCSCF, and response theory. Use as the reference text alongside
Module~III.2 from Lecture L07 onward.
\end{tcolorbox}
\end{bibentry}

\begin{bibentry}{t4light}
\textbf{Koch, W.\ \& Holthausen, M.~C.}\\
\textit{A Chemist's Guide to Density Functional Theory} (2nd ed.)\\
Wiley-VCH, 2001.\\[4pt]
\recommended\\[4pt]
\begin{tcolorbox}[annotbox]
Accessible DFT text for Tier~4 chemistry students. Good companion for
Module~III.3, Lectures L01--L05.
\end{tcolorbox}
\end{bibentry}

%------------------------------------------------------------
\subsection{Key Tier 4 Papers}
%------------------------------------------------------------

\begin{bibentry}{t4light}
\textbf{Hohenberg, P.\ \& Kohn, W.}\\
\textit{Inhomogeneous electron gas}\\
Phys.~Rev.~\textbf{136}, B864 (1964).\\[4pt]
\essential\ \ [\textbf{P}]\\[4pt]
\begin{tcolorbox}[annotbox]
Proves the two Hohenberg--Kohn theorems. Primary reading for Module~III.3 Lecture~L01.
\end{tcolorbox}
\end{bibentry}

\begin{bibentry}{t4light}
\textbf{Kohn, W.\ \& Sham, L.~J.}\\
\textit{Self-consistent equations including exchange and correlation effects}\\
Phys.~Rev.~\textbf{140}, A1133 (1965).\\[4pt]
\essential\ \ [\textbf{P}]\\[4pt]
\begin{tcolorbox}[annotbox]
The Kohn--Sham DFT paper. Primary reading for Module~III.3. Nobel lecture (1998)
provides accessible context. Essential for any Tier~4 quantum chemistry student.
\end{tcolorbox}
\end{bibentry}

\begin{bibentry}{t4light}
\textbf{Feigenbaum, M.~J.}\\
\textit{Quantitative universality for a class of nonlinear transformations}\\
J.~Stat.~Phys.~\textbf{19}, 25--52 (1978). [\textbf{P}]\\[4pt]
\essential\\[4pt]
\begin{tcolorbox}[annotbox]
Original discovery of $\delta\approx 4.6692$ and $\alpha$. Cross-reference for
Tier~4 nonlinear dynamics students. The QM programme analogue is Kohn \& Sham (1965).
\end{tcolorbox}
\end{bibentry}

\begin{bibentry}{t4light}
\textbf{Li, T.-Y.\ \& Yorke, J.~A.}\\
\textit{Period three implies chaos}\\
Am.~Math.~Monthly \textbf{82}, 985--992 (1975). [\textbf{P}]\\[4pt]
\essential\\[4pt]
\begin{tcolorbox}[annotbox]
Coined the word ``chaos''. Cross-reference Tier~4 example for illustrating how
short landmark papers can reshape a field (analogous to Born 1926 in QM).
\end{tcolorbox}
\end{bibentry}

\newpage
%#############################################################
\section{Tier 5 (PhD) --- Monographs, Reviews \& Primary Literature}
\label{sec:t5}
%#############################################################

\begin{tcolorbox}[tierheader=t5red]
Tier 5 --- PhD
\end{tcolorbox}

\medskip
\begin{tcolorbox}[keybox]
\textbf{Goal:} Achieve research-level mastery; read primary literature directly;
identify open problems; connect the programme to current research fronts.
Primary papers and monographs are the main reading; textbooks serve as reference.
Use alongside CPJG, WRQ, and ORQ artifact types from the PSM.
\end{tcolorbox}

%------------------------------------------------------------
\subsection{Functional Analysis \& Mathematical Foundations}
%------------------------------------------------------------

\begin{bibentry}{t5light}
\textbf{Reed, M.\ \& Simon, B.}\\
\textit{Methods of Modern Mathematical Physics} (4 vols.)\\
Academic Press, 1972--1979.\\[4pt]
\essential\ \ [\textbf{M}]\\[4pt]
\begin{tcolorbox}[annotbox]
Vol.~I (\textit{Functional Analysis}): Hilbert spaces, operators, spectral theory.
Vol.~II (\textit{Fourier Analysis, Self-Adjointness}): self-adjoint extensions,
deficiency indices, Gel'fand triple. The mathematical foundation for Module~I.2
Lecture~L12. Essential for Tier~5. Every PhD student should own these volumes.
\end{tcolorbox}
\end{bibentry}

\begin{bibentry}{t5light}
\textbf{Rudin, W.}\\
\textit{Functional Analysis} (2nd ed.)\\
McGraw-Hill, 1991.\\[4pt]
\essential\\[4pt]
\begin{tcolorbox}[annotbox]
Banach and Hilbert spaces, the spectral theorem for bounded operators. Chapter~12
(unbounded operators) is directly relevant to the self-adjoint extension problem
for $\hat{p}$, directly supporting Module~I.2 Lecture~L12 at PhD level.
\end{tcolorbox}
\end{bibentry}

\begin{bibentry}{t5light}
\textbf{Blank, J., Exner, P.\ \& Havl\'{\i}\v{c}ek, M.}\\
\textit{Hilbert Space Operators in Quantum Physics} (2nd ed.)\\
Springer, 2008.\\[4pt]
\essential\\[4pt]
\begin{tcolorbox}[annotbox]
The most physics-oriented rigorous treatment. Covers spectral measures, rigged
Hilbert spaces, and quantum dynamics at the PhD level. Directly supports
Module~I.2 Lectures L11--L12 and Module~I.3 Lectures L04--L08.
\end{tcolorbox}
\end{bibentry}

\begin{bibentry}{t5light}
\textbf{Gel'fand, I.~M.\ \& Vilenkin, N.~Ya.}\\
\textit{Generalized Functions, Vol.~4: Applications of Harmonic Analysis}\\
Academic Press, 1964.\\[4pt]
\recommended\ \ [\textbf{M}]\\[4pt]
\begin{tcolorbox}[annotbox]
Original source for the rigged Hilbert space (Gel'fand triple) framework.
Chapter~1 is the reference for Module~I.2 Lecture~L12 (``Rigged Hilbert Spaces
\& Mathematical Foundations'') at the PhD level.
\end{tcolorbox}
\end{bibentry}

%------------------------------------------------------------
\subsection{Advanced Quantum Mechanics}
%------------------------------------------------------------

\begin{bibentry}{t5light}
\textbf{Negele, J.~W.\ \& Orland, H.}\\
\textit{Quantum Many-Particle Systems}\\
Addison-Wesley, 1988 (Perseus reprint 1998).\\[4pt]
\essential\\[4pt]
\begin{tcolorbox}[annotbox]
Path integrals and Green's functions at graduate/research level. Supports
Modules~II.2 and~II.5.
\end{tcolorbox}
\end{bibentry}

\begin{bibentry}{t5light}
\textbf{Kadanoff, L.~P.\ \& Baym, G.}\\
\textit{Quantum Statistical Mechanics}\\
Benjamin, 1962 (Perseus reprint 1989).\\[4pt]
\essential\ \ [\textbf{M}]\\[4pt]
\begin{tcolorbox}[annotbox]
Original derivation of the Kadanoff--Baym equations. Primary reading for
Module~II.5 Lecture~L05.
\end{tcolorbox}
\end{bibentry}

\begin{bibentry}{t5light}
\textbf{Martin, R.~M., Reining, L.\ \& Ceperley, D.}\\
\textit{Interacting Electrons: Theory and Computational Approaches}\\
Cambridge University Press, 2016.\\[4pt]
\essential\\[4pt]
\begin{tcolorbox}[annotbox]
The definitive modern text connecting DFT, many-body perturbation theory, and
quantum Monte Carlo. Covers all of Series~III at research depth.
\end{tcolorbox}
\end{bibentry}

\begin{bibentry}{t5light}
\textbf{Haag, R.}\\
\textit{Local Quantum Physics: Fields, Particles, Algebras} (2nd ed.)\\
Springer, 1996.\\[4pt]
\recommended\ \ [\textbf{M}]\\[4pt]
\begin{tcolorbox}[annotbox]
Algebraic QFT; the deepest mathematical treatment of the measurement problem
and locality. Background for Series~II Module~II.4.
\end{tcolorbox}
\end{bibentry}

%------------------------------------------------------------
\subsection{Essential PhD-Level Papers (QM \& QC)}
%------------------------------------------------------------

\begin{bibentry}{t5light}
\textbf{Stone, M.~H.}\\
\textit{On one-parameter unitary groups in Hilbert space}\\
Ann.~Math.~\textbf{33}, 643--648 (1932).\\[4pt]
\essential\ \ [\textbf{P}]\\[4pt]
\begin{tcolorbox}[annotbox]
Stone's theorem on strongly continuous unitary groups. The mathematical basis
for $\hat{U}(t)=e^{-i\hat{H}t/\hbar}$. Primary reading for Module~I.3
Lecture~L04 at PhD level.
\end{tcolorbox}
\end{bibentry}

\begin{bibentry}{t5light}
\textbf{Gleason, A.~M.}\\
\textit{Measures on the closed subspaces of a Hilbert space}\\
J.~Math.~Mech.~\textbf{6}, 885--893 (1957).\\[4pt]
\essential\ \ [\textbf{P}]\\[4pt]
\begin{tcolorbox}[annotbox]
Derives the Born rule from measure theory on closed subspaces. The deepest
axiomatic justification of $P=|\langle a|\psi\rangle|^2$. Essential for
Module~I.2 Lecture~L07 at PhD level and directly relevant to the ORQ prompt
``What minimal operational axioms force complex Hilbert space?''
\end{tcolorbox}
\end{bibentry}

\begin{bibentry}{t5light}
\textbf{Lindblad, G.}\\
\textit{On the generators of quantum dynamical semigroups}\\
Commun.~Math.~Phys.~\textbf{48}, 119--130 (1976).\\[4pt]
\essential\ \ [\textbf{P}]\\[4pt]
\begin{tcolorbox}[annotbox]
Derives the Lindblad master equation from complete positivity. Primary reading
for Module~I.3 Lecture~L10 (``Open Systems and Decoherence'').
\end{tcolorbox}
\end{bibentry}

\begin{bibentry}{t5light}
\textbf{Kraus, K.}\\
\textit{States, Effects, and Operations: Fundamental Notions of Quantum Theory}\\
Springer Lecture Notes in Physics \textbf{190}, 1983.\\[4pt]
\recommended\ \ [\textbf{M}]\\[4pt]
\begin{tcolorbox}[annotbox]
Kraus operators, CPTP maps, and the operational formulation of quantum mechanics.
Background for Module~I.2 Lecture~L11 and Module~I.3 Lecture~L10.
\end{tcolorbox}
\end{bibentry}

%------------------------------------------------------------
\subsection{Nonlinear Dynamics PhD Cross-reference}
%------------------------------------------------------------

\begin{bibentry}{t5light}
\textbf{de Melo, W.\ \& van Strien, S.}\\
\textit{One-Dimensional Dynamics}\\
Springer, 1993. [\textbf{M}]\\[4pt]
\essential\\[4pt]
\begin{tcolorbox}[annotbox]
Definitive research monograph on interval maps. Proves ACIMs, renormalisation,
Feigenbaum theory. Exact analogue of Reed \& Simon for 1D dynamics.
\end{tcolorbox}
\end{bibentry}

\begin{bibentry}{t5light}
\textbf{Lyubich, M.}\\
\textit{Almost every real quadratic map is either regular or stochastic}\\
Ann.~Math.~\textbf{156}, 1--78 (2002). [\textbf{P}]\\[4pt]
\recommended\\[4pt]
\begin{tcolorbox}[annotbox]
Resolves the measure-theoretic dichotomy for the logistic family; one of the deepest
results in 1D dynamics.
\end{tcolorbox}
\end{bibentry}

\begin{bibentry}{t5light}
\textbf{Young, L.-S.}\\
\textit{What are SRB measures and which dynamical systems have them?}\\
J.~Stat.~Phys.~\textbf{108}, 733--754 (2002). [\textbf{R}]\\[4pt]
\essential\\[4pt]
\begin{tcolorbox}[annotbox]
Definitive survey of physical (SRB) invariant measures; directly analogous to
density operator surveys in QM.
\end{tcolorbox}
\end{bibentry}

\newpage
%#############################################################
\section{Full Alphabetical Catalogue}
\label{sec:fullcat}
%#############################################################

\begin{tcolorbox}[keybox]
Complete alphabetically ordered catalogue of all references across Tiers~1--5 and the
programme syllabus, grouped by category. Each entry carries its tier badge(s) and type tag.
\end{tcolorbox}

%------------------------------------------------------------
\subsection*{Category A \quad Popular Science \& Historical Narrative}
\addcontentsline{toc}{subsection}{Category A: Popular Science}

\begin{tcolorbox}[catbox]
\begin{enumerate}[leftmargin=2.5em, label=\textbf{[A\arabic*]}, itemsep=5pt]
  \item \tierbadge{t1green}{T1}\ Cox, B.\ \& Forshaw, J. \textit{The Quantum Universe.} Da Capo Press, 2011. [\textbf{N}]
  \item \tierbadge{t1green}{T1}\ Feynman, R.~P. \textit{QED: The Strange Theory of Light and Matter.} Princeton University Press, 1985. [\textbf{N}] \essential
  \item \tierbadge{t1green}{T1}\ Feynman, R.~P., Leighton, R.~B.\ \& Sands, M. \textit{The Feynman Lectures on Physics, Vol.~III.} Addison-Wesley, 1965. [\textbf{N}] \essential
  \item \tierbadge{t1green}{T1}\ Gleick, J. \textit{Chaos: Making a New Science.} Penguin Books, 1987. [\textbf{N}] \essential
  \item \tierbadge{t1green}{T1}\ Kumar, M. \textit{Quantum: Einstein, Bohr, and the Great Debate.} Norton, 2008. [\textbf{N}]
  \item \tierbadge{t1green}{T1}\ Penrose, R. \textit{The Emperor's New Mind.} Oxford University Press, 1989. [\textbf{N}]
  \item \tierbadge{t1green}{T1}\ Stewart, I. \textit{Does God Play Dice?} Blackwell, 1989. [\textbf{N}]
\end{enumerate}
\end{tcolorbox}

%------------------------------------------------------------
\subsection*{Category B \quad Undergraduate Textbooks (Tiers 1--3)}
\addcontentsline{toc}{subsection}{Category B: Undergraduate Textbooks}

\begin{tcolorbox}[catbox]
\begin{enumerate}[leftmargin=2.5em, label=\textbf{[B\arabic*]}, itemsep=5pt]
  \item \tierbadge{t2blue}{T2}\ Alligood, K.~T., Sauer, T.~D.\ \& Yorke, J.~A. \textit{Chaos: An Introduction to Dynamical Systems.} Springer, 1996.
  \item \tierbadge{t2blue}{T2}\ Binney, J.\ \& Skinner, D. \textit{The Physics of Quantum Mechanics.} Oxford University Press, 2014.
  \item \tierbadge{t2blue}{T2}\ Byron, F.~W.\ \& Fuller, R.~W. \textit{Mathematics of Classical and Quantum Physics.} Dover, 1992.
  \item \tierbadge{t3purple}{T3}\ Devaney, R.~L. \textit{An Introduction to Chaotic Dynamical Systems} (2nd ed.). Addison-Wesley, 1989.
  \item \tierbadge{t2blue}{T2}\ Griffiths, D.~J.\ \& Schroeter, D.~F. \textit{Introduction to Quantum Mechanics} (3rd ed.). Cambridge University Press, 2018. \essential
  \item \tierbadge{t3purple}{T3}\ Hannabuss, K. \textit{An Introduction to Quantum Theory.} Oxford University Press, 1997.
  \item \tierbadge{t3purple}{T3}\ Isham, C.~J. \textit{Lectures on Quantum Theory.} Imperial College Press, 1995.
  \item \tierbadge{t2blue}{T2}\ Jensen, F. \textit{Introduction to Computational Chemistry} (3rd ed.). Wiley, 2017.
  \item \tierbadge{t2blue}{T2}\ Shankar, R. \textit{Principles of Quantum Mechanics} (2nd ed.). Springer, 1994. \essential
  \item \tierbadge{t2blue}{T2}\ Strang, G. \textit{Introduction to Linear Algebra} (5th ed.). Wellesley--Cambridge Press, 2016. \essential
  \item \tierbadge{t2blue}{T2}\ Strogatz, S.~H. \textit{Nonlinear Dynamics and Chaos} (2nd ed.). Westview Press, 2015. \essential
  \item \tierbadge{t2blue}{T2}\ Townsend, J.~S. \textit{A Modern Approach to Quantum Mechanics} (2nd ed.). University Science Books, 2012.
\end{enumerate}
\end{tcolorbox}

%------------------------------------------------------------
\subsection*{Category C \quad Graduate Textbooks (Tiers 4--5)}
\addcontentsline{toc}{subsection}{Category C: Graduate Textbooks}

\begin{tcolorbox}[catbox]
\begin{enumerate}[leftmargin=2.5em, label=\textbf{[C\arabic*]}, itemsep=5pt]
  \item \tierbadge{t4orange}{T4}\ Ballentine, L.~E. \textit{Quantum Mechanics: A Modern Development} (2nd ed.). World Scientific, 2014.
  \item \tierbadge{t5red}{T5}\ Blank, J., Exner, P.\ \& Havl\'{\i}\v{c}ek, M. \textit{Hilbert Space Operators in Quantum Physics} (2nd ed.). Springer, 2008. \essential
  \item \tierbadge{t4orange}{T4}\ Bruus, H.\ \& Flensberg, K. \textit{Introduction to Many-Body Quantum Theory.} Oxford University Press, 2004. \essential
  \item \tierbadge{t4orange}{T4}\ Cohen-Tannoudji, C., Diu, B.\ \& Lalo\"e, F. \textit{Quantum Mechanics} (2 vols.). Wiley, 2019.
  \item \tierbadge{t4orange}{T4}\ Fetter, A.~L.\ \& Walecka, J.~D. \textit{Quantum Theory of Many-Particle Systems.} Dover, 2003. \essential
  \item \tierbadge{t4orange}{T4}\ Galindo, A.\ \& Pascual, P. \textit{Quantum Mechanics} (2 vols.). Springer, 1990--1991. \essential
  \item \tierbadge{t4orange}{T4}\ Helgaker, T., J\o rgensen, P.\ \& Olsen, J. \textit{Molecular Electronic-Structure Theory.} Wiley, 2000. \essential
  \item \tierbadge{t5red}{T5}\ Kadanoff, L.~P.\ \& Baym, G. \textit{Quantum Statistical Mechanics.} Perseus, 1989. [\textbf{M}] \essential
  \item \tierbadge{t4orange}{T4}\ Koch, W.\ \& Holthausen, M.~C. \textit{A Chemist's Guide to Density Functional Theory} (2nd ed.). Wiley-VCH, 2001.
  \item \tierbadge{t5red}{T5}\ Martin, R.~M., Reining, L.\ \& Ceperley, D. \textit{Interacting Electrons.} Cambridge University Press, 2016. \essential
  \item \tierbadge{t5red}{T5}\ Negele, J.~W.\ \& Orland, H. \textit{Quantum Many-Particle Systems.} Perseus, 1998. \essential
  \item \tierbadge{t5red}{T5}\ Rudin, W. \textit{Functional Analysis} (2nd ed.). McGraw-Hill, 1991. \essential
  \item \tierbadge{t4orange}{T4}\ Stefanucci, G.\ \& van Leeuwen, R. \textit{Nonequilibrium Many-Body Theory of Quantum Systems.} Cambridge University Press, 2013. \essential
  \item \tierbadge{t4orange}{T4}\ Szabo, A.\ \& Ostlund, N.~S. \textit{Modern Quantum Chemistry.} Dover, 1996. \essential
  \item \tierbadge{t4orange}{T4}\ Weinberg, S. \textit{Lectures on Quantum Mechanics} (2nd ed.). Cambridge University Press, 2015. \essential
\end{enumerate}
\end{tcolorbox}

%------------------------------------------------------------
\subsection*{Category D \quad Research Monographs (Tier 5 \& Research Track)}
\addcontentsline{toc}{subsection}{Category D: Research Monographs}

\begin{tcolorbox}[catbox]
\begin{enumerate}[leftmargin=2.5em, label=\textbf{[D\arabic*]}, itemsep=5pt]
  \item \tierbadge{t5red}{T5}\ Collet, P.\ \& Eckmann, J.-P. \textit{Iterated Maps on the Interval as Dynamical Systems.} Birkh\"auser, 1980 (reprint 2009). [\textbf{M}]
  \item \tierbadge{t5red}{T5}\ de Melo, W.\ \& van Strien, S. \textit{One-Dimensional Dynamics.} Springer, 1993. [\textbf{M}] \essential
  \item \tierbadge{t5red}{T5}\ Gel'fand, I.~M.\ \& Vilenkin, N.~Ya. \textit{Generalized Functions, Vol.~4.} Academic Press, 1964. [\textbf{M}]
  \item \tierbadge{t5red}{T5}\ Haag, R. \textit{Local Quantum Physics} (2nd ed.). Springer, 1996. [\textbf{M}]
  \item \tierbadge{t5red}{T5}\ Kraus, K. \textit{States, Effects, and Operations.} Springer, 1983. [\textbf{M}]
  \item \tierbadge{t5red}{T5}\ Reed, M.\ \& Simon, B. \textit{Methods of Modern Mathematical Physics} (4 vols.). Academic Press, 1972--1979. [\textbf{M}] \essential
  \item \tierbadge{t3purple}{T3}\ von Neumann, J. \textit{Mathematical Foundations of Quantum Mechanics.} Princeton University Press, 1955. [\textbf{M}]
\end{enumerate}
\end{tcolorbox}

%------------------------------------------------------------
\subsection*{Category E \quad Primary Research Papers}
\addcontentsline{toc}{subsection}{Category E: Primary Research Papers}

\begin{tcolorbox}[catbox]
\begin{enumerate}[leftmargin=2.5em, label=\textbf{[E\arabic*]}, itemsep=5pt]
  \item \tierbadge{t3purple}{T3}\ Born, M. Zur Quantenmechanik der Sto\ss vorg\"ange. \textit{Z.~Phys.} \textbf{37}, 863--867 (1926). [\textbf{P}] \essential
  \item \tierbadge{t5red}{T5}\ Eckmann, J.-P.\ \& Ruelle, D. Ergodic theory of chaos and strange attractors. \textit{Rev.~Mod.~Phys.} \textbf{57}, 617--656 (1985). [\textbf{R}]
  \item \tierbadge{t4orange}{T4}\ Feigenbaum, M.~J. Quantitative universality for a class of nonlinear transformations. \textit{J.~Stat.~Phys.} \textbf{19}, 25--52 (1978). [\textbf{P}] \essential
  \item \tierbadge{t5red}{T5}\ Gleason, A.~M. Measures on the closed subspaces of a Hilbert space. \textit{J.~Math.~Mech.} \textbf{6}, 885--893 (1957). [\textbf{P}] \essential
  \item \tierbadge{t4orange}{T4}\ Hohenberg, P.\ \& Kohn, W. Inhomogeneous electron gas. \textit{Phys.~Rev.} \textbf{136}, B864 (1964). [\textbf{P}] \essential
  \item \tierbadge{t4orange}{T4}\ Kohn, W.\ \& Sham, L.~J. Self-consistent equations including exchange and correlation. \textit{Phys.~Rev.} \textbf{140}, A1133 (1965). [\textbf{P}] \essential
  \item \tierbadge{t4orange}{T4}\ Li, T.-Y.\ \& Yorke, J.~A. Period three implies chaos. \textit{Am.~Math.~Monthly} \textbf{82}, 985--992 (1975). [\textbf{P}]
  \item \tierbadge{t5red}{T5}\ Lindblad, G. On the generators of quantum dynamical semigroups. \textit{Commun.~Math.~Phys.} \textbf{48}, 119--130 (1976). [\textbf{P}] \essential
  \item \tierbadge{t5red}{T5}\ Lyubich, M. Almost every real quadratic map is either regular or stochastic. \textit{Ann.~Math.} \textbf{156}, 1--78 (2002). [\textbf{P}]
  \item \tierbadge{t3purple}{T3}\ May, R.~M. Simple mathematical models with very complicated dynamics. \textit{Nature} \textbf{261}, 459--467 (1976). [\textbf{P}] \essential
  \item \tierbadge{t5red}{T5}\ Stone, M.~H. On one-parameter unitary groups in Hilbert space. \textit{Ann.~Math.} \textbf{33}, 643--648 (1932). [\textbf{P}] \essential
  \item \tierbadge{t3purple}{T3}\ Wigner, E.~P. On unitary representations of the inhomogeneous Lorentz group. \textit{Ann.~Math.} \textbf{40}, 149--204 (1939). [\textbf{P}]
  \item \tierbadge{t5red}{T5}\ Young, L.-S. Statistical properties of dynamical systems with some hyperbolicity. \textit{Ann.~Math.} \textbf{147}, 585--650 (1998). [\textbf{P}]
  \item \tierbadge{t5red}{T5}\ Young, L.-S. What are SRB measures and which dynamical systems have them? \textit{J.~Stat.~Phys.} \textbf{108}, 733--754 (2002). [\textbf{R}] \essential
\end{enumerate}
\end{tcolorbox}

%------------------------------------------------------------
\subsection*{Category F \quad Review Articles \& Surveys}
\addcontentsline{toc}{subsection}{Category F: Review Articles}

\begin{tcolorbox}[catbox]
\begin{enumerate}[leftmargin=2.5em, label=\textbf{[F\arabic*]}, itemsep=5pt]
  \item \tierbadge{t5red}{T5}\ Eckmann, J.-P. Roads to turbulence in dissipative dynamical systems. \textit{Rev.~Mod.~Phys.} \textbf{53}, 643--654 (1981). [\textbf{R}]
  \item \tierbadge{t5red}{T5}\ Lyubich, M. Forty years of unimodal dynamics. \textit{J.~Math.~Sci.} \textbf{170}, 680--713 (2010). [\textbf{R}]
  \item \tierbadge{t5red}{T5}\ Viana, M. Stochastic dynamics of deterministic systems. \textit{IMPA Lecture Notes} (1997). [\textbf{R}]
\end{enumerate}
\end{tcolorbox}

\bigskip
%------------------------------------------------------------
\subsection*{Tier--Category Quick-Index}
\addcontentsline{toc}{subsection}{Tier--Category Quick-Index}
%------------------------------------------------------------

\begin{center}
\begin{tabular}{lp{11cm}}
\toprule
\textbf{Tier} & \textbf{Recommended categories and key items}\\
\midrule
\tierbadge{t1green}{T1} HS
  & All of Cat.~A; Cat.~B: [B3, B5, B11]; no primary papers\\[3pt]
\tierbadge{t2blue}{T2} BegUG
  & Cat.~B: [B2, B5, B8, B9, B10, B11, B12]; Cat.~A as context;
    Cat.~E: [E10] (May 1976) abstract only\\[3pt]
\tierbadge{t3purple}{T3} AdvUG
  & Cat.~B (all); Cat.~C: [C4, C15]; Cat.~D: [D7];
    Cat.~E: [E1, E3, E10, E12] in full\\[3pt]
\tierbadge{t4orange}{T4} MSc
  & Cat.~C (all); Cat.~D: [D3, D6, D7];
    Cat.~E: [E1, E4, E5, E6, E8, E11, E12] as primary reading;
    Cat.~F: [F1, F2]\\[3pt]
\tierbadge{t5red}{T5} PhD
  & Cat.~D (all); Cat.~E (complete); Cat.~F (all);
    Cat.~C as reference; current arXiv preprints\\[3pt]
Research track
  & Cat.~D (all); Cat.~E (all); Cat.~F (all);
    journal club selections; current arXiv\\
\bottomrule
\end{tabular}
\end{center}

%=============================================================
\end{document}
